% bab2.tex — gabungan BAB_II_A/B/C *_POLISHED.md (## A + ## B + ## C)

\section*{BAB II --- PEMBAHASAN}
\addcontentsline{toc}{section}{BAB II --- PEMBAHASAN}

\subsection*{A. Periodisasi Sejarah Perkembangan Hadis}

\subsubsection*{1. Periodisasi Masa Nabi Muhammad SAW (610--632 M / 13 SH--11 H)}

Ciri khas penyaluran hadis pada era kenabian terletak pada keragaman jalurnya: hafalan, pengamalan ibadah, \textit{qira'ah}, \textit{sama'}, dan pencatatan perdana berjalan beriringan sebagai satu ekosistem penjagaan (Tanjung \& Fadhlurrahman, 2025; Musadad, 2026). Hadis bukan sekadar ujaran tekstual, melainkan \textit{living sunnah} yang dihidupkan lewat suri teladan Nabi dalam salat, muamalah, dan pergaulan sosial (Razak et al., 2025; Tanjung \& Fadhlurrahman, 2025). Kelisanan yang kukuh pada masyarakat Arab yang \textit{ummi} menjadikan ingatan kolektif sebagai sarana utama, sedangkan tulisan sekadar menyangga memori, bukan kodifikasi sistematis (Musadad, 2026; Mosa, 2025). Kewibawaan bertumpu pada kehadiran ragawi Nabi selaku sumber hidup, sehingga rantai pembuktian formal belum dituntut (Tanjung \& Fadhlurrahman, 2025; Razak et al., 2025). Pewarisan hadis menempuh empat jalur: \textit{sama'} (menyimak lalu menghafal), \textit{qira'ah}/\textit{'ard} (melafalkan kembali di hadapan Nabi untuk dikoreksi), pengajaran interaktif melalui majelis dan khotbah, serta penyaksian partisipatif dalam praktik wudu, salat, dan haji (Tanjung \& Fadhlurrahman, 2025; Razak et al., 2025) (Musadad, 2026; Mosa, 2025). Corak interaktif-dialogis itu, dipadukan dengan ekologi lisan-tulis yang saling melengkapi, menegaskan bahwa pemisahan kaku antara oral dan tertulis tidaklah memadai (Razak et al., 2025; Tanjung \& Fadhlurrahman, 2025) (Tanjung \& Fadhlurrahman, 2025; Musadad, 2026). Fase Nabi karenanya tampil sebagai tahapan formatif yang mutlak, mewariskan memori, praktik, dan catatan yang terserak sebagai bahan mentah bagi seluruh penilaian angkatan berikutnya (Razak et al., 2025; Mosa, 2025) (Musadad, 2026; Tanjung \& Fadhlurrahman, 2025).

Persoalan penulisan hadis diwarnai tarik-menarik dua klaster riwayat: pelarangan mencatat selain Al-Qur'an berhadapan dengan perkenan tegas untuk mendokumentasikan sabda Nabi (Tanjung \& Fadhlurrahman, 2025; Razak et al., 2025). Pelarangan awal dimaknai sebagai \textit{sadd al-dzari'ah}, upaya menghindari bercampurnya wahyu yang masih turun bertahap dengan perkataan Nabi di tengah umat yang baru mengenal aksara (Tanjung \& Fadhlurrahman, 2025; Musadad, 2026). Seiring menguatnya keberaksaraan Madinah dan kesanggupan membedakan Al-Qur'an dari hadis, Nabi memberikan kelonggaran bersyarat (\textit{rukhsah}), sehingga larangan dan perkenan dipahami sebagai strategi berjenjang, bukan pertentangan (Razak et al., 2025; Tanjung \& Fadhlurrahman, 2025). Kompromi ini sejalan dengan penilaian bahwa tradisi tulis perdana memang telah ada, meski mutu keaslian tiap naskah bertingkat dan tidak seluruhnya dapat direkonstruksi secara kukuh (Musadad, 2026; Razak et al., 2025). Klaim yang meniadakan penulisan hadis pada abad pertama Hijriah karenanya terlalu jauh, sebab cikal dokumentasi tertulis telah tumbuh berdampingan dengan kelisanan (Tanjung \& Fadhlurrahman, 2025; Mosa, 2025).

Kalangan sahabat terbagi dalam dua peran fungsional yang saling melengkapi: penulis dan penghafal (Tanjung \& Fadhlurrahman, 2025; Musadad, 2026). Abdullah bin Amr bin al-Ash tampil sebagai penulis utama melalui \textit{Sahifah al-Sadiqah} yang direstui Nabi, bersama Piagam Madinah dan surat-surat diplomatik yang menegaskan keberaksaraan komunitas perdana (Tanjung \& Fadhlurrahman, 2025; Razak et al., 2025). Para penghafal seperti Abu Hurairah, Aisyah, Anas bin Malik, dan Ibnu Abbas bersandar pada kekuatan memori dan ketekunan bergaul dengan Nabi, sehingga menjadi titik sentral transmisi (Mosa, 2025; Musadad, 2026). Keragaman logat kesukuan dan keterbatasan ingatan menumbuhkan periwayatan \textit{bi al-ma'na} di samping \textit{bi al-lafdzi} --- kemestian sosiolinguistik yang kelak menuntut kritik matan (Razak et al., 2025; Rozikin, 2026). Wafat Nabi pada 11 H/632 M mengakhiri otoritas tunggal yang sentralistis dan membuka era sahabat yang tercerai-berai, sehingga pembuktian dan pewilayahan transmisi menjadi agenda pokok berikutnya (Tanjung \& Fadhlurrahman, 2025; Liew, 2025).

\subsubsection*{2. Periodisasi Masa Sahabat (11--80 H / 632--700 M)}

Watak menonjol era sahabat ialah penyebaran geografis sepeninggal Nabi yang mengubah penyaluran sentralistis-Madinah menjadi pola kedaerahan menuju Syam, Irak, Mesir, dan Persia (Tanjung \& Fadhlurrahman, 2025; Liew, 2025). Para sahabat menetap di kota-kota baru sebagai pengajar, hakim, dan gubernur, sehingga tiap kawasan menumbuhkan khazanah dan kecenderungan pemahaman yang khas (Liew, 2025; Mosa, 2025). Implikasinya, lahir pewarisan bercorak regional berikut cikal \textit{rihlah}: para pemburu hadis menempuh perjalanan menjumpai sahabat tertentu demi riwayat yang otoritatif (Tanjung \& Fadhlurrahman, 2025; Razak et al., 2025). Kedekatan pribadi dan martabat \textit{suhbah} berubah menjadi modal kewibawaan pokok, sebab hadis belum dihimpun dalam kitab besar (Musadad, 2026; Mosa, 2025). Kian beragamnya penyaluran memupuk tuntutan tata cara pembuktian yang berkesadaran metode (Razak et al., 2025; Tanjung \& Fadhlurrahman, 2025). Penyebaran tersebut memadat pada empat titik utama --- Madinah yang konservatif sebagai pewaris amalan Nabi, Makkah yang menonjol lewat kewibawaan Ibnu Abbas, serta Kufah dan Basrah sebagai simpul gesit tempat riwayat dan persoalan hukum perkotaan bertemu, hingga perdebatan \textit{ra'yu} kelak memanas (Tanjung \& Fadhlurrahman, 2025; Mosa, 2025) (Musadad, 2026; Liew, 2025) (Liew, 2025; Razak et al., 2025). Keterhubungan antarpusat ditopang perjalanan ilmu yang kian gencar (Tanjung \& Fadhlurrahman, 2025; Mosa, 2025). Titik-titik kewilayahan ini menjelma warisan struktural yang segera diwarisi tabi'in, manakala jalinan personal yang mulai renggang beralih menjadi kelembagaan guru-murid dengan kaidah \textit{isnad} yang kian mengeras (Musadad, 2026; Razak et al., 2025).

Tuntutan pembuktian dijawab lewat \textit{tathabbut} para Khulafaur Rasyidin dengan ragam prosedur yang bermuara pada tujuan tunggal: menjaga keutuhan informasi (Tanjung \& Fadhlurrahman, 2025; Razak et al., 2025). Abu Bakar mensyaratkan kesaksian penyerta (\textit{bayyinah}) sebelum menerima riwayat asing, sebagai wujud tanggung jawab merawat kemurnian ajaran sepeninggal Nabi (Razak et al., 2025; Liew, 2025). Umar bahkan menolak dan memperingatkan pembawa hadis yang janggal tanpa penguat, sehingga tercatat sebagai tonggak skeptisisme metodologis terhadap pernyataan perseorangan (Tanjung \& Fadhlurrahman, 2025; Mosa, 2025). Ali menempuh cara penyumpahan (\textit{yamin}): perawi yang belum dikenal diminta bersumpah atas nama Allah sebagai ujian kejujuran berdimensi psikologis-moral (Razak et al., 2025; Tanjung \& Fadhlurrahman, 2025). Rangkaian praktik itu secara tegas merupakan cikal kritik \textit{sanad} --- nalar kritis angkatan pertama atas keterpercayaan sumber, meski belum dirumuskan sebagai disiplin yang terstruktur (Mosa, 2025; Razak et al., 2025).

Faktor pengakselerasi pelembagaan pembuktian ialah \textit{al-fitnah al-kubra} pascaterbunuhnya Usman, yang memecah umat ke dalam kubu politik-teologis yang saling berebut pengesahan (Liew, 2025; Tanjung \& Fadhlurrahman, 2025). Perseteruan Ali--Muawiyah dan kelompok sempalan garis keras menyulut pemalsuan (\textit{wad'}): kisah direkayasa lalu disandarkan kepada Nabi demi menopang klaim sektarian (Liew, 2025; Razak et al., 2025). \textit{Afterlife} hadis khilafah tiga puluh tahun membuktikan bahwa satu \textit{matan} dapat diproduksi ulang untuk menopang kisah kekuasaan yang saling bertolak belakang, sehingga krisis keaslian berwatak struktural (Liew, 2025; Musadad, 2026). Sebagai benteng pertahanan, sahabat dan tabi'in perdana mulai menuntut penyebutan mata rantai penyaluran --- peralihan dari keyakinan mutlak menuju keraguan metodologis (Tanjung \& Fadhlurrahman, 2025; Mosa, 2025). Jejaknya terekam dalam pernyataan Ibnu Sirin: pemeriksaan \textit{sanad} diberlakukan sesudah fitnah, lengkap dengan kaidah penerimaan berbasis afiliasi --- diambil dari Ahlussunnah dan ditampik dari ahli bidah (Liew, 2025; Razak et al., 2025).

\subsubsection*{3. Periodisasi Masa Tabi'in (80--130 H / 700--750 M)}

Watak era tabi'in ialah peralihan dari regionalisasi menuju jejaring keilmuan lintas kawasan yang direkatkan oleh \textit{rihlah fi talab al-'ilm} (Tanjung \& Fadhlurrahman, 2025; Mosa, 2025). Kewibawaan berpindah dari martabat \textit{suhbah} menuju keterandalan jalinan penyaluran guru-ke-murid lintas kota (Musadad, 2026; Razak et al., 2025). Hijaz menonjol dalam memelihara amalan Madinah, sedangkan Irak tampil dalam wacana hukum-teologi yang tanggap terhadap persoalan perkotaan (Tanjung \& Fadhlurrahman, 2025; Liew, 2025). Kegesitan para pencari ilmu meluaskan penyaluran sekaligus mempertemukan ragam riwayat lintas jalur dalam satu forum (Mosa, 2025; Razak et al., 2025). Era tabi'in dengan demikian menandai tahapan pelembagaan: hadis beranjak dari pusaka personal-komunal menjadi kegiatan keilmuan yang terorganisasi (Musadad, 2026; Tanjung \& Fadhlurrahman, 2025). Pelembagaan ini disangga tiga wadah: \textit{halaqah} masjid, majelis kediaman ulama, dan jalinan guru-murid melalui \textit{sama'}, \textit{qira'ah}, \textit{ijazah} (Tanjung \& Fadhlurrahman, 2025; Musadad, 2026). Dalam \textit{halaqah}, syekh membacakan riwayat, murid menyimak, mencatat, lalu mencocokkannya lewat pembacaan ulang demi pengabsahan (Musadad, 2026; Razak et al., 2025). Majelis khusus memungkinkan penyaluran yang lebih intensif dan selektif, termasuk \textit{ijazah} sebagai lisensi meriwayatkan (Tanjung \& Fadhlurrahman, 2025; Mosa, 2025). Kekerabatan lintas angkatan melahirkan silsilah terdokumentasi, sehingga \textit{sanad} mencerminkan ikatan pedagogis nyata yang terverifikasi secara historis (Mosa, 2025; Musadad, 2026). Yang tertulis dan yang lisan bekerja berdampingan: catatan murid tetap memerlukan pengabsahan lisan dari guru demi kewibawaan penuh (Musadad, 2026; Tanjung \& Fadhlurrahman, 2025).

Seiring melebarnya jejaring, \textit{isnad} mengeras dari tuntutan yang longgar menjadi patokan teknis pembuktian, disertai tunas \textit{'ilm al-rijal} (Razak et al., 2025; Mosa, 2025). Ulama tabi'in mencatat siapa menerima riwayat dari siapa, di mana perjumpaan itu berlangsung, serta mutu hafalan dan integritas tiap mata rantai sebagai arsip penilaian (Mosa, 2025; Tanjung \& Fadhlurrahman, 2025). Peneguhan ini tak terlepas dari pertikaian sektarian dan derasnya pernyataan riwayat pasca-fitnah yang menuntut atribusi ketat pemisah Ahlussunnah dari ahli bidah (Liew, 2025; Razak et al., 2025). Pembandingan ragam \textit{matan} lintas jalur riwayat dijalankan untuk mengendus penyimpangan redaksional --- cikal kritik \textit{sanad-cum-matan} mutakhir (Razak et al., 2025; Al-Imam, 2026). Akibatnya, \textit{isnad} tumbuh sebagai jawaban rangkap, ilmiah sekaligus sosiologis, dalam merawat keutuhan teks di tengah perebutan kewibawaan yang tajam (Tanjung \& Fadhlurrahman, 2025; Mosa, 2025).

Puncak pelembagaan ialah \textit{tadwin} resmi yang disokong negara pada era Umar bin Abdul Aziz (99--101 H/717--720 M), yang risau atas wafatnya para penghafal dan maraknya riwayat yang lemah (Tanjung \& Fadhlurrahman, 2025; Razak et al., 2025). Titah resmi kepada para gubernur --- terutama Abu Bakar bin Hazm di Madinah --- dan kepada al-Zuhri memerintahkan penghimpunan, penyaringan, dan pembukuan hadis yang terserak (Tanjung \& Fadhlurrahman, 2025; Musadad, 2026). Karya perdana umumnya masih berupa \textit{jami'}/\textit{musannaf} bersahaja yang merangkum hadis \textit{marfu'}, fatwa sahabat (\textit{mauquf}), dan pandangan tabi'in (\textit{maqtu'}) ke dalam bab fikih --- belum terhimpun sebagai kompilasi profetik murni seperti abad ketiga (Musadad, 2026; Razak et al., 2025). \textit{al-Muwatta'} karya Malik memadukan hadis dan amal Madinah, sehingga menunjukkan gerak kewibawaan amalan setempat menuju perekaman tekstual lintas kawasan (Musadad, 2026; Tanjung \& Fadhlurrahman, 2025). \textit{Tadwin} tampil sebagai titik balik dari yang oral-personal menuju yang skriptural-komunal, sekaligus fondasi bagi ledakan kanonisasi pada era ulama klasik (Tanjung \& Fadhlurrahman, 2025; Razak et al., 2025).

\subsubsection*{4. Periodisasi Masa Ulama Klasik (130--235 H / 750--850 M)}

Watak era ulama klasik ialah kanonisasi: titik kewibawaan berpindah dari perawi perseorangan menuju kitab-kitab kanonik beserta metode kritik yang mapan (Musadad, 2026; Razak et al., 2025). Penghimpunan menyeleksi ratusan ribu riwayat yang beredar melalui asas pengabsahan yang tegas dan terdokumentasi (Tanjung \& Fadhlurrahman, 2025; Al-Imam, 2026). Ragam penulisan berkembang dari \textit{musannaf} campuran menuju himpunan murni (\textit{mujarrad}) yang memilahkan sabda Nabi dari fatwa sahabat dan opini tabi'in, diawali kitab \textit{musnad} seperti \textit{Musnad} Ahmad yang disusun berdasar nama sahabat (Musadad, 2026; Razak et al., 2025). Puncaknya, \textit{Kutub al-Sittah} --- Bukhari, Muslim, Abu Dawud, Tirmidzi, Nasa'i, Ibnu Majah --- kukuh sebagai rujukan mutlak Sunni hingga era mutakhir (Al-Imam, 2026; Rozikin, 2026). Kewibawaan dilembagakan dalam artefak tekstual terstandar yang dapat digandakan, diajarkan, dan diuji lintas angkatan (Musadad, 2026; Tanjung \& Fadhlurrahman, 2025).

Para tokoh utama mempersonifikasikan patokan tersebut (Musadad, 2026; Razak et al., 2025). Al-Bukhari (w. 256 H/870 M) menyaring perawi secara ketat dengan mendayagunakan sumber lisan-tertulis melalui \textit{sama'}, \textit{ijazah}, \textit{wajadah}, \textit{'ard}, sehingga keutuhan dan kecermatan menjadi syarat yang tak dapat ditawar (Tanjung \& Fadhlurrahman, 2025; Mosa, 2025). Muslim (w. 261 H/875 M) menonjol dalam penyandingan lintas jalur riwayat demi keruntutan redaksi, sedangkan Abu Dawud, Nasa'i, dan Ibnu Majah meluaskan kanon ke hukum praktis melalui \textit{sunan} (Razak et al., 2025; Rozikin, 2026). Tirmidzi menyumbangkan taksonomi lewat penamaan \textit{hasan} sebagai kategori tengah antara sahih dan \textit{dhaif}, mengatasi kebuntuan dikotomi biner (Al-Imam, 2026; Razak et al., 2025). Polemik \textit{afterlife} hadis khilafah tiga puluh tahun membuktikan bahwa kewibawaan para imam kanonik turut menstabilkan ingatan kolektif di tengah perebutan tafsir politik (Liew, 2025; Musadad, 2026).

Metodologi penyangganya ialah perpaduan kritik \textit{sanad} dan \textit{matan} melalui \textit{al-jarh wa al-ta'dil} (Razak et al., 2025; Al-Imam, 2026). Kritik \textit{sanad} menakar kesinambungan mata rantai, \textit{'adalah}, dan \textit{dabt}; kritik \textit{matan} menakar keselarasan terhadap Al-Qur'an, riwayat yang lebih kukuh, akal, dan fakta kesejarahan (Al-Imam, 2026; Rozikin, 2026). Telaah bandingan kritik \textit{matan} dalam \textit{Fath al-Bari} (Ibn Hajar) dan \textit{Majallat al-Manar} (Rasyid Rida) menegaskan bahwa kepedulian terhadap isi berakar pada khazanah klasik, bukan reka baru zaman modern --- polarisasi tradisionis-reformis hanya mempertajam rumusannya pada era modern (Al-Imam, 2026; Razak et al., 2025). Penyandingan aneka jalur \textit{sanad} dan ragam \textit{matan} ala Muslim membuktikan bahwa kritik klasik berwatak dialektis, tidak berhenti pada nama perawi (Razak et al., 2025; Mosa, 2025). Anggapan bahwa ulama klasik semata-mata peduli \textit{sanad} dengan demikian merupakan salah kaprah: keduanya saling melengkapi sejak tahap kematangan (Al-Imam, 2026; Tanjung \& Fadhlurrahman, 2025). Terobosan epistemologisnya berupa trikotomi sahih-\textit{hasan}-\textit{dhaif} dengan \textit{Sahihayn} sebagai patokan emas (Razak et al., 2025; Al-Imam, 2026). Prasyarat penerimaan --- kesinambungan \textit{sanad}, keutuhan dan kecermatan perawi, serta bebas dari \textit{syudzudz} dan \textit{'illah} --- mengalihkan pengabsahan yang intuitif menjadi tata cara intersubjektif yang dapat direplikasi (Razak et al., 2025; Rozikin, 2026). Kanonisasi \textit{Sahihayn} turut ditentukan oleh arsitektur kitab yang memudahkan pemakaian dan penerimaan komunitas lintas angkatan (Musadad, 2026; Tanjung \& Fadhlurrahman, 2025). Daya tahan taksonomi ini terbukti dalam wacana waris mutakhir, yang kesahihan dalilnya masih ditentukan oleh kategori abad ketiga Hijriah (Rozikin, 2026; Al-Imam, 2026). Kemapanan itu kemudian membeku menjadi skolastik repetitif, sehingga modernitas --- skeptisisme Barat dan disrupsi digital --- menuntut pembaruan paradigma (Liew, 2025; Matin Bin Salman, 2025).

\subsubsection*{5. Periodisasi Masa Modern (2015--2026)}

Watak era modern, yang terekam dalam literatur 2015--2026, ialah pluralisasi metode yang memadukan kritik \textit{sanad}, \textit{matan}, pembacaan kontekstual, dan ikhtiar digital-komputasional yang interdisipliner (Razak et al., 2025; Matin Bin Salman, 2025). Perkembangan kritik teks memetakan peralihan dari dominasi \textit{sanad} menuju \textit{sanad-cum-matan} yang peka terhadap \textit{asbab al-wurud} dan kemaslahatan (Razak et al., 2025; Al-Imam, 2026). Hadis dimengerti tidak sekadar sebagai teks pewarisan, melainkan sebagai wacana yang diproduksi dan dirundingkan dalam ekosistem yang terdesentralisasi (Matin Bin Salman, 2025; Razak et al., 2025). Kritik rangkap atas \textit{sanad}-\textit{matan} menentukan kesahihan dalil legislasi mutakhir, misalnya dalam fikih waris (Rozikin, 2026; Al-Imam, 2026). Era modern meluaskan khazanah klasik: etos \textit{tathabbut}, \textit{rihlah}, dan \textit{jarh-ta'dil} dialihbahasakan menjadi perangkat hermeneutik-komputasional baru (Razak et al., 2025; Mosa, 2025).

Rintangan bersumber pada dua gelombang: skeptisisme orientalis dan guncangan digital (Al-Imam, 2026; Razak et al., 2025). Gelombang pertama --- Goldziher dan Schacht melalui teori \textit{projecting back}, yang dimatangkan Juynboll lewat \textit{common link} yang mencurigai ruas di atas periwayat poros (Razak et al., 2025; Mosa, 2025) --- dijawab dengan pembelaan empiris atas kesinambungan \textit{tadwin} abad pertama serta \textit{isnad-cum-matan} yang mengaitkan percabangan \textit{isnad} dengan ragam \textit{matan} lintas geografis (Razak et al., 2025; Al-Imam, 2026). Perbandingan tradisionis (\textit{Fath al-Bari}) versus reformis (\textit{Majallat al-Manar}) mengenai hadis gejala alam memperlihatkan polarisasi internal: mempertahankan kanon berhadapan dengan penafsiran ulang yang rasional (Al-Imam, 2026; Razak et al., 2025). Perdebatan kaum skeptis dan apologetis justru memacu perumusan kerangka yang lebih canggih ketimbang skolastik yang beku (Razak et al., 2025; Matin Bin Salman, 2025). Gelombang kedua --- mediatisasi lewat media sosial dan AI yang menggeser \textit{sanad} otoritatif menuju kewenangan algoritmik (Matin Bin Salman, 2025; Alghamdi et al., 2025) --- menghadirkan peredaran infografis dan video ringkas tanpa \textit{takhrij}, sehingga viralitas tanpa keabsahan menggoyang kewibawaan tradisional (Matin Bin Salman, 2025; Razak et al., 2025). Jawabannya bersifat komputasional: model \textit{pretrained} pendeteksi hadis palsu dengan ketepatan bermakna, berpijak pada ciri kebahasaan-struktural (Alghamdi et al., 2025; Mosa, 2025), ditambah graf pengetahuan-\textit{transformer} untuk disambiguasi narator dan pengendus anomali \textit{sanad} otomatis dalam skala luas (Mosa, 2025; Alghamdi et al., 2025). Meski begitu, komputasi tetap memerlukan pengabsahan ulama, sebab mesin probabilistik tidak memahami aspek syariat-kesejarahan secara kognitif (Alghamdi et al., 2025; Razak et al., 2025).

Sintesis kajian ini menunjukkan bahwa studi hadis melaju menuju \textit{living hadith} dan hermeneutika digital yang menempatkan teknologi sebagai lanjutan etos klasik (Matin Bin Salman, 2025; Razak et al., 2025). \textit{Living hadith} membedah hadis sebagai nilai hidup dalam laku sosial --- dari penafsiran ulang hadis misoginis hingga teologi sebagai daya tahan psikologis (Matin Bin Salman, 2025; Rozikin, 2026). Hermeneutika digital menuntut literasi rangkap: filologis klasik (\textit{takhrij}, kritik \textit{sanad-matan}) sekaligus verifikasi digital atas sumber siber dan deteksi pemalsuan algoritmik (Alghamdi et al., 2025; Mosa, 2025). Evolusi kritik teks, artikulasi kewenangan profetik dalam \textit{al-Muwatta'} dan \textit{musannaf}, serta peran \textit{tadwin} menegaskan benang merah ketahanan metodologis: dari \textit{sahifah} hingga pangkalan data digital (Musadad, 2026; Tanjung \& Fadhlurrahman, 2025; Razak et al., 2025). Pemetaan lima periode ini menjadi landasan analitis bagi uraian penghafalan dan penghimpunan hadis, dari kelisanan hingga kodifikasi resmi (Tanjung \& Fadhlurrahman, 2025; Razak et al., 2025).


\subsection*{B. Penghafalan dan Penghimpunan Hadis}

\subsubsection*{1. Tradisi Menghafal Hadis (\textit{Hifz})}

Masyarakat Arab pra-Islam meletakkan legitimasi ilmu pada kelisanan: syair disimak, dituturkan, dan dihafal jauh sebelum aksara menyebar meluas (Rashwan, 2024; Jaiyeoba \& Osmani, 2024). Praktik itu diwarisi Islam perdana melalui majelis, khotbah, dan perjumpaan langsung dengan Nabi, sehingga sabda, tindakan, dan ketetapan beliau dijaga terutama lewat \textit{hifz} (Jaiyeoba \& Osmani, 2024; Watukila, 2024). Penyampaian hadis pada generasi awal bersifat performatif: lafal disertai isyarat tubuh, intonasi suara, dan konteks didaktis yang menghidupkan teks (Rashwan, 2024; Gil, 2024). Dokumentasi tertulis sebenarnya sudah dikenal, tetapi otoritasnya tetap tersubordinasi pada budaya simak dan koreksi personal; primasi lisan tidak tergantikan (Lutfianto et al., 2024; Jaiyeoba \& Osmani, 2024). Dengan demikian, \textit{hifz} menjadi fondasi kultural bagi transmisi ilmu agama lintas generasi, melampaui sekadar teknik menghafal individual (Rashwan, 2024; Watukila, 2024). Tradisi ini ditopang oleh perangkat pedagogis baku --- \textit{sama'}, \textit{qira'ah} atau \textit{'ard}, \textit{tasmi'}, \textit{tikrar}, dan \textit{musalsal} yang ritualistik (Jaiyeoba \& Osmani, 2024; Rashwan, 2024). Dalam \textit{sama'} murid menyerap materi yang didengarnya lewat repetisi; dalam \textit{qira'ah} murid menyetor bacaannya untuk dikoreksi dan diotorisasi guru sebelum berhak meriwayatkan (Jaiyeoba \& Osmani, 2024; Gil, 2024). Repetisi \textit{tasmi'} mengonsolidasikan hafalan, sedangkan \textit{musalsal} --- pewarisan adab penyampaian sekaligus isi --- menunjukkan bahwa metode periwayatan turut terekam dalam memori transmisi (Rashwan, 2024; Watukila, 2024). Verifikasi diperkuat oleh lembar \textit{sahifah}, bukti simak, dan kolasi: hafalan disandingkan dengan tulisan sebelum murid dipandang layak mengantongi \textit{ijazah} (Lutfianto et al., 2024; Jaiyeoba \& Osmani, 2024). Singkatnya, hafalan pada masa awal merupakan kerja kolektif yang diregulasi oleh jaringan guru--murid dan majelis, bukan olah individual yang soliter (Gil, 2024; Rashwan, 2024).

Kelebihan \textit{hifz} terletak pada daya lentur dan vitalitasnya: hadis tetap lestari sebagai tradisi hidup yang kontekstual-didaktis, terus berdenyut kendati infrastruktur tulis masih minim (Rashwan, 2024; Jaiyeoba \& Osmani, 2024). Sifat performatifnya memungkinkan adaptasi dialektal melalui \textit{bi al-ma'na} --- kelenturan sosiolinguistik yang masuk akal mengingat audiensnya lintas kabilah (Rashwan, 2024; Lutfianto et al., 2024). Sebaliknya, hafalan juga rentan bervariasi: materi dapat terinterpolasi, susunan lafal tertukar, atau nisbah periwayat keliru di sepanjang rantai transmisi (Jaiyeoba \& Osmani, 2024; Watukila, 2024). Ingatan lisan mengandung risiko bias sistematis; tanpa kontrol, hafalan semata tidak menjamin otentisitas jangka panjang (Rashwan, 2024; Obeid \& Hussein, 2023). Inilah paradoks \textit{hifz}: lincah sebagai medium hidup, namun tetap memerlukan topangan arsip tertulis serta kritik perawi agar validitasnya teruji (Jaiyeoba \& Osmani, 2024; Lutfianto et al., 2024).

Justru dari kerapuhan itulah \textit{hifz} menyumbang pada mekanisme verifikasi: defisit memori dikompensasi oleh \textit{sanad} sebagai teknologi sosial penjamin mutu hafalan tiap perawi (Watukila, 2024; Jaiyeoba \& Osmani, 2024). Jaringan \textit{sanad} memiliki dwifungsi --- validasi dan pemetaan sosial: ia merangkai guru--murid, kota-kota ilmu, dan strata generasi, sekaligus mengaudit otoritas setiap riwayat (Gil, 2024; Watukila, 2024). Mutu hafalan dikaitkan dengan kritik perawi melalui dwisyarat \textit{'adalah} dan \textit{dabt}; perawi yang lemah ingatannya tersisih setelah disandingkan dengan yang lebih kokoh (Obeid \& Hussein, 2023; Abbas \& Rawabdeh, 2025). Seleksi hafalan semacam ini terbukti menopang kanonisasi al-Bukhari, yang menyeleksi secara ketat akurasi estafet lisan-tulis (Obeid \& Hussein, 2023; Jaiyeoba \& Osmani, 2024). Alhasil, tradisi hafalan bukan rival kritik ilmiah, melainkan rahim kelahiran \textit{al-jarh wa al-ta'dil}; dari basis oral inilah wacana menyambung ke dokumentasi era Nabi (Rashwan, 2024; Watukila, 2024).

\subsubsection*{2. Penulisan Hadis pada Masa Nabi (\textit{Kitabah})}

Komunitas Arab pada era Nabi memiliki keberaksaraan yang terbatas: hanya sebagian kecil sahabat yang melek baca-tulis di tengah mayoritas yang bertutur lisan (Lutfianto et al., 2024; Watukila, 2024). Tulisan berfungsi sebagai alat mnemonik --- nota ringkas, \textit{sahifah} individual, antologi fragmen --- yang mengikat materi seiring jaringan periwayatan terus berekspansi (Lutfianto et al., 2024; Jaiyeoba \& Osmani, 2024). Migrasi dari hafalan ke tulisan berlangsung secara bertahap, dengan \textit{qira'ah}, \textit{sama'}, dan sesi simak sebagai penghubungnya (Lutfianto et al., 2024; Rashwan, 2024). Historiografi \textit{mutaqaddimin} versi Mustafa Azami menegaskan bahwa penulisan berlangsung kontinu tanpa henti sejak abad pertama Hijriah (Watukila, 2024; Lutfianto et al., 2024). Artinya, \textit{kitabah} pada masa kenabian melekat dalam ekologi \textit{hifz}, bukan menjadi antitesisnya: amalan \textit{kitabah} menyangga, memantik ingatan, dan mengawal sirkulasi lisan (Jaiyeoba \& Osmani, 2024; Rashwan, 2024).

Persoalan intinya terletak pada dua rumpun riwayat yang tampak kontradiktif: larangan menulis selain Al-Qur'an, berhadapan dengan otorisasi eksplisit untuk mencatat sabda Nabi (Lutfianto et al., 2024; Watukila, 2024). Larangan yang paling valid dinisbahkan kepada Abu Sa'id al-Khudri, sedangkan otorisasi dinisbahkan kepada Abdullah bin Amr, yang mengklaim lisannya tidak pernah mengucapkan selain kebenaran (Lutfianto et al., 2024; Jaiyeoba \& Osmani, 2024). Al-Khatib al-Baghdadi, dalam \textit{Taqyid al-'Ilm}, mengumpulkan kedua kelompok riwayat tersebut dan menyimpulkan bahwa hanya satu riwayat larangan yang valid, berbanding riwayat pembolehan yang jauh lebih banyak (Watukila, 2024; Lutfianto et al., 2024). Sebagian kecil orientalis memanfaatkan riwayat larangan itu untuk menyimpulkan bahwa hadis baru dibukukan pada abad ketiga --- kesimpulan yang mengabaikan korpus riwayat otorisasi dan bukti \textit{sahifah} (Watukila, 2024; Rashwan, 2024). Yang dipertaruhkan bukan ada-tidaknya kontroversi itu sendiri, melainkan bagaimana kedua rumpun riwayat yang tampak berseberangan tersebut dibaca secara kritis, historis, dan metodologis (Lutfianto et al., 2024; Jaiyeoba \& Osmani, 2024). Rekonsiliasi yang disepakati mayoritas ulama memposisikan larangan tersebut sebagai ketentuan yang bersifat kontekstual dan partikular, bukan veto total atas aktivitas dokumentasi (Lutfianto et al., 2024; Watukila, 2024). Larangan itu menjawab tiga kekhawatiran konkret: kerancuan antara hadis dan wahyu yang sedang turun, teralihkannya perhatian dari primasi wahyu, dan tercampurnya kedua teks pada satu medium (Watukila, 2024; Jaiyeoba \& Osmani, 2024). Kuncinya bersifat kronologis: larangan berlaku pada fase awal wahyu, ketika ayat-ayat yang ringkas masih sulit dibedakan dari hadis, sedangkan otorisasi terbit pada fase akhir, ketika sebagian besar wahyu telah turun dan umat mampu mengenali idiom masing-masing (Lutfianto et al., 2024; Watukila, 2024). Kaidahnya: sirnanya \textit{'illat} meruntuhkan interdiksi, sehingga norma general yang \textit{mujmal} dinasakh oleh otorisasi partikular bagi sahabat yang presisi (Watukila, 2024; Lutfianto et al., 2024). Begitu preservasi Al-Qur'an mencapai titik optimal, \textit{ijma'} melegitimasi penulisan dan kodifikasi hadis secara terbuka dan luas tanpa keraguan (Jaiyeoba \& Osmani, 2024; Watukila, 2024).

Bukti-bukti awal dari berbagai sumber memperkuat rekonsiliasi ini; yang terkuat adalah \textit{Sahifah al-Sadiqah} milik Abdullah bin Amr, yang ditulis atas izin tegas dari Nabi (Lutfianto et al., 2024; Watukila, 2024). Korpus serupa milik Ali, Jabir, dan sejumlah sahabat lain juga tercatat, disertai arsip diplomasi --- Piagam Madinah dan korespondensi kepada raja-raja asing --- yang menjadi saksi keberaksaraan komunitas mula-mula (Watukila, 2024; Jaiyeoba \& Osmani, 2024). Sahabat tidak berhenti pada menghafal: mereka juga membukukan, mempekerjakan penulis, menghadirkan saksi, dan mencatatkan tarikh --- profil penghafal ternyata tidak identik dengan keengganan terhadap aksara (Lutfianto et al., 2024; Rashwan, 2024). Fragmen-fragmen itu beredar lintas generasi sebagai bahan primer bagi antologi kanonik, dengan mata rantai tekstual yang tetap dapat dilacak (Obeid \& Hussein, 2023; Watukila, 2024). Arsip awal ini menyediakan basis empiris bagi kajian penghimpunan pada generasi sahabat, saat nota-nota personal ikut berpindah bersama pemiliknya (Gil, 2024; Lutfianto et al., 2024).

\subsubsection*{3. Penghimpunan Hadis pada Masa Sahabat}

Determinan utama penghimpunan pada era sahabat adalah eksodus ke luar Madinah: hafalan dan tulisan turut bermigrasi, melahirkan himpunan-himpunan kewilayahan yang bersifat desentralistis (Gil, 2024; Watukila, 2024). Ekspansi ke Syam, Irak, Mesir, dan Persia menumbuhkan korpus lokal yang bertumpu pada figur sahabat yang bermukim sebagai pengajar, hakim, maupun gubernur --- ikatan yang bersifat personal dan renggang, bukan sentralisasi pembukuan (Gil, 2024; Kara, 2026). Genre \textit{buldan} mendokumentasikan \textit{Rihlah}, yaitu mobilitas penuntut ilmu antarkota menuju sahabat-sahabat kunci --- prototipe perjalanan keilmuan sekaligus agregasi transregional (Gil, 2024; Jaiyeoba \& Osmani, 2024). Karena energi komunitas lebih tersedot pada pembukuan Al-Qur'an, hadis terakumulasi melalui \textit{sahifah} individual, diktat murid, dan hafalan komunal tanpa patronase negara (Lutfianto et al., 2024; Watukila, 2024). Dispersi geografis ini memicu divergensi periwayatan, sehingga filtrasi yang ketat menjadi keharusan untuk memurnikan ajaran dari anomali (Jaiyeoba \& Osmani, 2024; Gil, 2024). Tabi'in menerima estafet berupa nota personal dan memori kolektif sebagai bahan mentah yang kemudian diteruskan, dibandingkan, dan ditata secara sistematis (Gil, 2024; Lutfianto et al., 2024). Kanal utamanya adalah kedekatan guru--murid antara sahabat dan tabi'in, sebagaimana termaktub dalam karya-karya proto seperti \textit{atraf} yang menjadi bekal sebelum majelis (Gil, 2024; Jaiyeoba \& Osmani, 2024). Prosopografi lingkaran Basrah membuktikan bahwa transmisi pada masa transisi ini lebih difasilitasi oleh tabi'in yang solid ketimbang sahabat lokal, dengan \textit{rihlah} mikro antarsentra yang memperkaya simpul-simpulnya (Gil, 2024; Watukila, 2024). Lembaran \textit{sahifah}, bukti simak, dan fragmen sahabat dikurasi dalam \textit{halaqah} tabi'in menuju tekstualisasi yang bersifat transregional (Lutfianto et al., 2024; Jaiyeoba \& Osmani, 2024). Estafet berlapis ini membuka jalan bagi kodifikasi resmi, tatkala penguasa era Umar bin Abdul Aziz menghimpun materi yang berserak ke dalam program yang lebih terencana (Watukila, 2024; Lutfianto et al., 2024).

Filtrasi ini dikenal dengan istilah \textit{tathabbut}, yaitu etos kehati-hatian Khulafaur Rasyidin yang mensyaratkan bukti tambahan sebelum mengadopsi kabar dari luar (Watukila, 2024; Jaiyeoba \& Osmani, 2024). Abu Bakar mensyaratkan \textit{bayyinah} berupa kesaksian tambahan; Umar mendiskualifikasi kabar ganjil yang tidak didukung penguat, bahkan hingga taraf ancaman; sementara Ali mensyaratkan \textit{yamin} bagi perawi yang tidak dikenal (Watukila, 2024; Lutfianto et al., 2024). Kebijakan diskualifikasi tersebut merintis skeptisisme yang metodis, karena testimoni tunggal mesti lolos uji intersubjektif sebelum dapat diterima sebagai hujah yang sah (Jaiyeoba \& Osmani, 2024; Abbas \& Rawabdeh, 2025). Perluasannya dapat dipetakan pada doktrin \textit{munkar} dalam kritik al-Dzahabi terhadap al-Hakim di \textit{al-Mustadrak}: riwayat yang berlawanan dengan perawi yang lebih superior gugur, sekalipun rantai eksternalnya tetap kontinu (Abbas \& Rawabdeh, 2025; Obeid \& Hussein, 2023). Dengan demikian, \textit{tathabbut} merupakan prototipe kritik \textit{sanad} yang definitif, kendati idiom teknisnya belum baku pada masa itu (Watukila, 2024; Gil, 2024).

Sahabat senior berfungsi sentral sebagai simpul yang menyambungkan memori kenabian dengan generasi murid tabi'in yang datang kemudian (Gil, 2024; Kara, 2026). Aisyah tampil sebagai otoritas korektif-hermeneutis, yang menyanggah, mengklarifikasi, dan memaknai ulang kesaksian koleganya dengan modal pengalaman langsung dan konsistensi makna; artinya, kritik pemahaman internal telah berdenyut sejak generasi awal (Kara, 2026; Ardig et al., 2025). Ibnu Abbas unggul di Makkah dalam bidang tafsir-fikih, Anas memiliki memori khidmah yang ekstensif, dan Abu Hurairah tampil paling prolifik dengan korpus terluas, meskipun arsitektur kuantitatif jaringannya masih menuntut kajian prosopografi lanjutan (Gil, 2024; Watukila, 2024). Kontroversi seputar sakralitas Madinah membuktikan bahwa otoritas senior mengonsolidasi memori kolektif di tengah rivalitas tafsir antarregion yang tajam (Kara, 2026; Gil, 2024). Glosa dan koreksi mereka menjadi prototipe syarah yang pada periode berikutnya berkembang menjadi genre otonom dengan tradisi tersendiri (Ardig et al., 2025; Kara, 2026).

\subsubsection*{4. Kodifikasi Hadis Secara Resmi (\textit{Tadwin})}

Kelahiran \textit{tadwin} dipicu oleh dua kondisi darurat sekaligus: kekhawatiran punahnya khazanah setelah gugurnya \textit{huffaz} di medan perang, dan kontaminasi hadis apokrif-politis pasca-\textit{fitnah} (Lutfianto et al., 2024; Watukila, 2024). Menjelang akhir abad pertama, para saksi langsung mulai memudar seiring bermekarannya klaim-klaim sektarian; preservasi melalui memori privat dan arsip yang tercecer pun terbukti kurang memadai (Watukila, 2024; Kara, 2026). Pada era sahabat, pembukuan belum menjadi tuntutan karena fokus masih tersedot pada Al-Qur'an; barulah pada angkatan tabi'in dokumentasi menjadi mendesak guna membendung fabrikasi (Lutfianto et al., 2024; Jaiyeoba \& Osmani, 2024). Inilah momen infleksi: verifikasi privat-informal model \textit{tathabbut} harus dieskalasi menjadi agenda kolektif yang terorganisasi (Jaiyeoba \& Osmani, 2024; Lutfianto et al., 2024). Kondisi darurat tersebut memicu munculnya respons negara sebagai solusi institusional atas krisis preservasi yang cakupannya kian meluas (Watukila, 2024; Obeid \& Hussein, 2023). Titah Umar bin Abdul Aziz (99--101 H/717--720 M) dialamatkan kepada para gubernur --- terutama Abu Bakar bin Hazm di Madinah --- serta kepada al-Zuhri, dengan perintah untuk mengumpulkan, menyeleksi, dan membukukan hadis yang tercecer (Lutfianto et al., 2024; Watukila, 2024). Surat khalifah tersebut mengamanatkan penyisiran sistematis atas hafalan dan diktat ulama lintas region --- sebuah intervensi berpatronase negara yang menggeser inisiatif privat menjadi proyek komunal (Watukila, 2024; Jaiyeoba \& Osmani, 2024). Al-Zuhri, sebagai eksekutor, mengagregasi bahan antarkota dan merumuskan fondasi metodologisnya, meski ia tidak layak direduksi menjadi poros tunggal kodifikasi dalam episode panjang yang mendahuluinya (Lutfianto et al., 2024; Gil, 2024). \textit{Rihlah} dioptimalkan sebagai wahana ekstraksi, dengan delegasi yang menyambangi repositori hidup di sentra-sentra ilmu yang berpencar (Gil, 2024; Lutfianto et al., 2024). Titah tersebut menjadi prasasti institusional yang menggeser \textit{Tadwin} dari sekadar diskursus menjadi kebijakan negara, sekaligus mematenkan \textit{tadwin} sebagai agenda komunal (Watukila, 2024; Jaiyeoba \& Osmani, 2024).

Produk abad kedua masih bersifat transisional: antologi \textit{jami'} dan \textit{musannaf} berformat bab-bab fikih, belum sepenuhnya steril sebagai kumpulan sabda Nabi semata (Lutfianto et al., 2024; Obeid \& Hussein, 2023). Ketiga kategori --- \textit{marfu'}, \textit{mauquf}, dan \textit{maqtu'} --- masih terbaur, sehingga otoritas profetik belum terdemarkasi dari fikih angkatan awal yang mengapitnya (Obeid \& Hussein, 2023; Ardig et al., 2025). \textit{al-Muwatta'} karya Malik diproyeksikan sebagai manifes konsensus amal Madinah, sementara \textit{Musannaf} karya 'Abd al-Razzaq berekspansi menuju agregasi yang lebih komprehensif, tekstual, dan transregional (Obeid \& Hussein, 2023; Lutfianto et al., 2024). Korpus transisional ini memasok bahan mentah bagi kodifikasi kanonik abad ketiga: al-Bukhari mendayagunakan sumber lisan-tulis melalui \textit{sama'}, \textit{ijazah}, \textit{wajadah}, dan \textit{'ard} (Obeid \& Hussein, 2023; Abbas \& Rawabdeh, 2025). Amalgamasi kategorial ini, secara ironis, justru membidani kebutuhan akan genre taksonomis baru yang memisahkan graduasi otentisitas secara eksplisit (Ardig et al., 2025; Obeid \& Hussein, 2023).

Konsekuensi historisnya bersifat paradigmatik: supremasi oral-privat bermutasi menjadi rezim skriptural-komunal garapan negara bersama komunitas ilmu (Watukila, 2024; Jaiyeoba \& Osmani, 2024). Otoritas bergeser dari kedekatan personal dengan perawi menuju kitab yang direplikasi, disirkulasikan kurikulumnya lintas region, dan diuji melalui kritik \textit{sanad} dan \textit{matan} secara formal (Obeid \& Hussein, 2023; Abbas \& Rawabdeh, 2025). Parameter autentikasi --- kontinuitas \textit{sanad}, integritas dan presisi perawi, serta sterilitas dari \textit{syudzudz} dan \textit{'illah} --- memuncak dalam kanonisasi \textit{Kutub al-Sittah} dan purifikasi \textit{mujarrad} (Obeid \& Hussein, 2023; Jaiyeoba \& Osmani, 2024). Kanonisasi ini sekaligus mengonsolidasikan memori kolektif dan mendemarkasi ortodoksi di tengah divergensi transmisi peninggalan \textit{tadwin} (Kara, 2026; Ardig et al., 2025). Pascakodifikasi, wacana bermigrasi menuju krisis suksesi: kemapanan arsip tertulis justru menginkubasi fabrikasi mutakhir, yang kemudian menuntut tipologi pemalsuan beserta desain penyelamatannya (Abbas \& Rawabdeh, 2025; Obeid \& Hussein, 2023).


\subsection*{C. Timbulnya Pemalsuan Hadis dan Upaya Penyelamatan}

\subsubsection*{1. Pengertian Hadis \textit{Maudu'}}

Kodifikasi formal mengubah transmisi lisan yang tadinya terserak menjadi korpus terdokumentasi, tetapi konsolidasi ini justru membuka celah bagi infiltrasi artifisial sehingga batasan \textit{maudu'} perlu ditegaskan (Muhamed Ali et al., 2015; Wasman et al., 2023). Hadis \textit{maudu'} dimaknai sebagai pemberitaan fabrikan tanpa basis historis yang dinisbatkan secara dusta kepada Nabi, padahal beliau tidak pernah menyabdakannya (Muhamed Ali et al., 2015; Calgan, 2024). Batasan tersebut menegaskan dua unsur pembentuknya: konstruksi \textit{matan} artifisial dan klaim nasab palsu melalui \textit{sanad} rekayasa yang dikesankan tersambung (Wasman et al., 2023; Mohamed \& Sarwar, 2022). Berbeda dari kekhilafan yang tak disengaja, \textit{maudu'} mensyaratkan unsur kesengajaan, sehingga penilaiannya mempersoalkan integritas etik transmiter, bukan sekadar daya ingatnya (Muhamed Ali et al., 2015; Shaaban et al., 2026). \textit{Maudu'} karenanya menempati kategori mandiri dalam taksonomi hadis, bukan sekadar tingkatan terendah dari hadis lemah (Wasman et al., 2023; Noor, 2023).

Kedudukan \textit{maudu'} menjadi lebih jelas ketika dikontraskan dengan \textit{dhaif}, \textit{matruk}, dan \textit{munkar}, yang masing-masing memuat gradasi kecacatan berbeda (Muhamed Ali et al., 2015; Wasman et al., 2023). Kategori \textit{dhaif} bersumber dari defek \textit{sanad} --- baik berupa keterputusan rantai maupun rapuhnya hafalan --- meskipun status ini terkadang dapat terangkat melalui penguat lintas jalur atau sokongan yurisprudensial (Noor, 2023; Tarmom et al., 2022). Sebaliknya, \textit{matruk} dan \textit{munkar} merepresentasikan kerusakan yang lebih serius: transmiternya dituduh berdusta, atau materinya bertentangan dengan riwayat perawi yang jauh lebih superior, sejalan dengan telaah \textit{munkar} atas autentikasi al-Hakim dalam \textit{al-Mustadrak} sebagai indikasi disfungsi \textit{isnad} yang berat (Wasman et al., 2023; Muhamed Ali et al., 2015). Berbeda dari ketiganya, \textit{maudu'} merupakan kreasi deliberatif yang dihimpun tersendiri dalam genre khusus (\textit{kutub al-maudu'at}), sebab terbukti dirancang secara sengaja, bukan sekadar lemah (Noor, 2023; Calgan, 2024). Pembedaan ini membawa implikasi komputasional: model deteksi otomatis harus mampu memisahkan kerapuhan biasa dari signature linguistik hasil rekayasa (Azmi et al., 2019; Mohamed \& Sarwar, 2022).

Konsekuensinya bersifat imperatif: konsensus ulama (\textit{ijma'}) mengharamkan secara mutlak transmisi \textit{maudu'} tanpa penjelasan terbuka mengenai statusnya kepada pendengar (Muhamed Ali et al., 2015; Wasman et al., 2023). Larangan ini berpijak pada ancaman Nabi berstatus \textit{mutawatir} bahwa pendusta atas namanya telah menyiapkan tempat duduknya di neraka --- formulasi yang merefleksikan jawaban reaktif atas praktik fabrikasi yang telah berlangsung (Wasman et al., 2023; Noor, 2023). Telaah kontekstual menegaskan bahwa kedudukan suatu riwayat menentukan legitimasi pemakaiannya sebagai dasar hukum, materi dakwah, maupun landasan keutamaan amal (Wasman et al., 2023; Calgan, 2024). Telaah Ibn Kasir atas sebagian materi \textit{Sahihayn} pada segmen \textit{sirah} membuktikan bahwa korpus kanonik pun tetap tunduk pada evaluasi substansi, sehingga toleransi pemakaian \textit{maudu'} demi \textit{fada'il al-a'mal} tetap ditolak oleh mayoritas ulama (Calgan, 2024; Noor, 2023). Ketegasan definisi, taksonomi, dan implikasi ini menghantarkan pembahasan pada problem historis: momentum kemunculan awal fabrikasi yang terorganisasi (Wasman et al., 2023; Su, 2021). Kriteria identifikasi klasik terkristalisasi sebagai jawaban atas gelombang pemalsuan yang aktual --- telaah transmiter 'Ali b. al-Madini menguat ketika klaim pasca-fitnah melimpah, sehingga menuntut diferensiasi metodis antara riwayat otentik dan fabrikan (Su, 2021; Noor, 2023; Muhamed Ali et al., 2015). Uraian kemudian bertransisi dari sisi definisional ke historis: silsilah fabrikasi sejak era kenabian hingga kritik kontemporer (Wasman et al., 2023; Su, 2021).

\subsubsection*{2. Sejarah Munculnya Pemalsuan Hadis}

Debat mengenai titik awal fabrikasi mempertemukan dua posisi: segelintir kalangan menempatkan embrionya pada masa Nabi masih hidup, sementara mayoritas memosisikannya pasca-fitnah (Wasman et al., 2023; Su, 2021). Kelompok minoritas ini bersandar pada peringatan Nabi terhadap para pendusta sebagai indikasi sosiologis bahwa fabrikasi duniawi pernah berlangsung, atau tengah dicoba (Wasman et al., 2023; Muhamed Ali et al., 2015). Namun, kajian filologis-historis yang rigor menunjukkan bahwa laporan spesifik tentang fabrikasi di era Nabi mengandung cacat \textit{sanad} fatal serta pelaku yang tidak jelas, sehingga tidak cukup kuat dijadikan evidensi bagi fabrikasi yang sistematis (Su, 2021; Noor, 2023). Mayoritas menegaskan bahwa ekosistem sejak era Nabi hingga awal masa Umar relatif steril dari fabrikasi massif, berkat kedekatan temporal dengan wahyu, kontrol internal, dan soliditas moral kolektif sahabat (Wasman et al., 2023; Muhamed Ali et al., 2015). Kajian ini mengikuti pandangan mayoritas tersebut: pemalsuan insidental pada era Nabi tidak setara dengan industri terstruktur yang muncul belakangan (Su, 2021; Wasman et al., 2023).

Titik demarkasinya ialah peristiwa \textit{al-fitnah al-kubra}: terbunuhnya Usman dan perang saudara antara Ali dan Muawiyah (Wasman et al., 2023; Su, 2021). Kekosongan otoritas dan krisis keabsahan yang menyertainya memadatkan kelompok-kelompok radikal --- cikal Syiah fanatik pendukung Ali dan Khawarij yang menentang keduanya --- yang membutuhkan fondasi teologis bagi narasi politik-militer mereka (Wasman et al., 2023; Noor, 2023). Ketika legitimasi dari Qur'an maupun riwayat sahih tidak ditemukan, frustrasi doktrinal ini mendorong rekayasa hadis untuk menjustifikasi posisi sendiri sekaligus meruntuhkan otoritas lawan (Wasman et al., 2023; Muhamed Ali et al., 2015). Irak --- khususnya Kufah dan Basrah --- bertransformasi menjadi sentra produksi terbesar, sehingga otoritas semisal Malik mencurigai ekstremitas transmisi dari Irak dan mewajibkan verifikasi bertingkat (Su, 2021; Noor, 2023). Ali sendiri, di Siffin, memperingatkan pasukannya tentang pencideraan integritas sabda Nabi (Wasman et al., 2023; Su, 2021).

Memasuki abad kedua hingga ketiga, seiring konsolidasi Umayyah dan transisi menuju Abbasiyah, fabrikasi meluas dari kontestasi elektoral menjadi industri lintas mazhab dan sekte --- Murji'ah, Qadariyah, Jabariyah --- serta persaingan hegemoni antarmazhab hukum (Wasman et al., 2023; Noor, 2023). Produksi yang tak terkendali ini tersalurkan melalui para \textit{qussas}, ambisi popularitas dan keuangan, serta pasar devosi seputar keutamaan amal, mukjizat, hingga eskatologi (Wasman et al., 2023; Muhamed Ali et al., 2015). Untuk meredam lonjakan atribusi pada abad kedua--ketiga ini, standardisasi kritik \textit{'ilal} dan \textit{rijal} dirintis oleh 'Ali b. al-Madini (Su, 2021; Muhamed Ali et al., 2015). Abad kedua hingga ketiga dengan demikian merepresentasikan transfigurasi pemalsuan dari instrumen faksi menjadi industri kultural yang meresap ke semua lapisan umat Muslim (Wasman et al., 2023; Noor, 2023). Linimasa ini --- dari embrio insidental melalui politik fitnah hingga industri lintas-sekte --- menegaskan bahwa fabrikasi merupakan simptom krisis otoritas, bukan insiden personal (Wasman et al., 2023; Alam \& Schneider, 2020). Pemahaman historis tersebut menjadi prasyarat bagi pembedahan motif, sebab determinan pendorongnya hanya dapat terbaca dalam konteks krisis yang menghasilkannya (Wasman et al., 2023; Su, 2021). Uraian selanjutnya bergeser dari pertanyaan kapan ke mengapa: klasifikasi motif struktural di balik produksi dan sirkulasi riwayat palsu (Muhamed Ali et al., 2015; Wasman et al., 2023).

\subsubsection*{3. Faktor-Faktor yang Mendorong Pemalsuan}

Variabel politik menempati posisi paling determinan: hadis, sebagai otoritas puncak setelah Al-Qur'an, menghimpun legitimasi besar bagi perebutan kekuasaan (Wasman et al., 2023; Muhamed Ali et al., 2015). Faksi-faksi yang memperebutkan keimamahan, khilafah, dan keutamaan khalifah merancang riwayat tandingan guna menopang kedudukan sekaligus mendiskreditkan lawan --- hadis difungsikan sebagai alat perang propaganda (Wasman et al., 2023; Su, 2021). Semakin pivotal peran hadis sebagai penentu identitas politik-keagamaan, semakin kuat pula dorongan untuk mencetak tandingannya; nalar pasar otoritas inilah yang menerangkan derasnya fabrikasi di Kufah, Basrah, Wasit, dan Baghdad (Wasman et al., 2023; Alam \& Schneider, 2020). Analisis jaringan transmiter \textit{Sahih Bukhari} membuktikan dimensi strukturalnya: pola sentralisasi dan percabangan jalur merefleksikan perebutan pengaruh di lingkaran politik-keilmuan (Alam \& Schneider, 2020; Kamran et al., 2023). Fabrikasi politik karenanya merupakan taktik kolektif kubu-kubu yang memperebutkan otoritas normatif, bukan penyimpangan akhlak perseorangan (Wasman et al., 2023; Muhamed Ali et al., 2015).

Variabel teologis-sektarian memperluas arena politik menjadi konflik doktrinal, meliputi isu \textit{qadar}, kehendak bebas, kedudukan pelaku dosa besar, hingga supremasi antarmazhab fikih (Wasman et al., 2023; Noor, 2023). Kontroversi seputar \textit{qadar} dan iradah dinyatakan secara eksplisit turut memupuk rekayasa dan manipulasi periwayatan, sebab tiap kubu membutuhkan sokongan tekstual profetik (Wasman et al., 2023; Muhamed Ali et al., 2015). Rivalitas antara Sunni, Syiah, dan Khawarij beserta turunannya mendorong penciptaan riwayat pemuliaan tokoh, pelaknatan terhadap seteru, dan klaim eskatologis yang mengokohkan identitas (Wasman et al., 2023; Su, 2021). Telaah kontekstual menegaskan bahwa narasi polemis semacam ini hanya dapat dievaluasi secara proporsional melalui \textit{asbab al-wurud}, bukan melalui literalisme a-historis (Wasman et al., 2023; Calgan, 2024). Kebangkitan tradisi Mamluk melalui \textit{Muqaddimah} karya Ibn al-Salah memperlihatkan bahwa kontroversi warisan senantiasa menuntut pembakuan ulang, karena akumulasi sektarian terus diwariskan antargenerasi (Noor, 2023; Wasman et al., 2023). Determinan destruktif-infiltratif melengkapi gambaran ini: unsur \textit{zindiq} beserta kalangan heretik menginfiltrasi komunitas untuk merusak kewibawaan arus utama dari dalam (Wasman et al., 2023; Muhamed Ali et al., 2015). Sebagian fabrikasi dikaitkan dengan lingkungan menyimpang, termasuk kemungkinan unsur \textit{zindiq} yang menyisipkan laporan penggerus kredibilitas syariat, walaupun saluran operasionalnya tidak selalu terdokumentasi secara rinci (Wasman et al., 2023; Noor, 2023). Berbeda dari kalangan sektarian yang masih mengimani kelompoknya sendiri, fabrikator destruktif berambisi mengaburkan Islam secara keseluruhan, sebab muatannya berkontradiksi dengan prinsip \textit{kulli} syariat (Wasman et al., 2023; Calgan, 2024). Kritik \textit{matan} Ibn Kasir atas segmen \textit{sirah} dalam \textit{Sahihayn} relevan untuk mendeteksi infiltrasi semacam ini: laporan yang lolos dari seleksi \textit{sanad} tetap dapat digugurkan melalui pengujian koherensi historis-teologis (Calgan, 2024; Wasman et al., 2023). Pemodelan gaya profetik berbasis \textit{AraBERT} menyediakan instrumen modern untuk mendeteksi penyimpangan stilistik infiltratif secara statistik (Shaaban et al., 2026; Mohamed \& Sarwar, 2022).

Variabel keempat, yang bercorak sosio-ekonomi-devosional --- mencakup motif dakwah yang keliru melalui \textit{targhib-tarhib}, kepentingan material, dan kebutuhan akan literatur populer --- justru memikul tanggung jawab terbesar atas bertahannya sirkulasi \textit{maudu'} hingga sekarang (Wasman et al., 2023; Muhamed Ali et al., 2015). Fabrikator kerap berdalih konstruktif: laporan palsu tentang keutamaan amal dinilai bermanfaat untuk mendorong kesalehan, sehingga kreasi demi dakwah yang mulia dianggap terlegitimasi (Wasman et al., 2023; Noor, 2023). Para \textit{qussas}, forum penceritaan, dan dakwah populer berperan sebagai saluran distribusi yang efektif dalam memenuhi hasrat moral-emosional-eskatologis di luar supervisi ahli hadis (Wasman et al., 2023; Muhamed Ali et al., 2015). Motivasi material memperburuk keadaan: nama Nabi dimanfaatkan untuk mendulang popularitas dan sokongan finansial dari penguasa maupun audiens (Muhamed Ali et al., 2015; Wasman et al., 2023). Kombinasi antara motif ideologis di pihak produsen dan motif populer di pihak sirkulator menjelaskan mengapa kemapanan kritik tidak serta-merta memberantas \textit{maudu'} --- pembahasan selanjutnya berlanjut pada ciri-ciri hadis palsu (Wasman et al., 2023; Azmi et al., 2019).

\subsubsection*{4. Ciri-Ciri Hadis Palsu}

Kriteria \textit{sanad} bertumpu pada reliabilitas transmiter, kontinuitas mata rantai, dan pendeteksian cacat (\textit{'illah}) melalui disiplin \textit{al-jarh wa al-ta'dil} (Muhamed Ali et al., 2015; Su, 2021). Indikator terkuatnya meliputi konfesi eksplisit dari pelaku, transmiter yang dikenal pendusta atau pemalsu, pretensi transmisi langsung yang berseberangan dengan catatan biografis, diskontinuitas rantai, dan transmiter \textit{majhul} (Muhamed Ali et al., 2015; Wasman et al., 2023). Praktik pemeriksaan tanggal lahir dan wafat, inventarisasi hubungan guru-murid, serta verifikasi \textit{liqa'} sejak generasi awal telah dibuktikan oleh 'Ali b. al-Madini, yang menggugurkan klaim perjumpaan yang mustahil secara kronologis (Su, 2021; Muhamed Ali et al., 2015). Analisis jaringan transmiter \textit{Sahih Bukhari} memodernkan pendekatan ini: anomali relasional --- baik berupa keterisolasian dari gugus seangkatan maupun trayektori yang deviatif dari pola kolektif --- dibaca sebagai isyarat statistik adanya distorsi (Alam \& Schneider, 2020; Kamran et al., 2023). Ontologi \textit{SemanticHadith} memformalkan pendekatan tersebut: transmiter, karya, dan jalur direpresentasikan sebagai jaringan semantik sehingga anomalinya dapat dikalkulasi (Kamran et al., 2023; Azmi et al., 2019).

Kriteria \textit{matan} menguji substansi riwayat terhadap Al-Qur'an, hadis sahih, fakta sejarah, akal, konsensus ulama, dan prinsip umum syariat (Wasman et al., 2023; Calgan, 2024). Gejala kewaspadaannya meliputi: absennya riwayat dalam koleksi yang diakui, redaksi yang janggal atau berbahasa rusak, makna yang mustahil secara ontologis, kontradiksi dengan evidensi yang lebih kuat, konten sektarian atau tidak senonoh, serta imbalan pahala atau ancaman yang tidak proporsional untuk amal yang sepele (Wasman et al., 2023; Mohamed \& Sarwar, 2022). Kritik Ibn Kasir atas segmen \textit{sirah} dalam \textit{Sahihayn} membuktikan bahwa kekuatan \textit{sanad} tidak selalu final: rantai yang kukuh pun dapat gugur melalui pengujian koherensi historis (Calgan, 2024; Wasman et al., 2023). Evaluasi fitur linguistik melalui \textit{machine learning} menegaskan relevansi kontemporer pendekatan ini: sintaksis, diversitas kosakata, dan deviasi stilistik terbukti berkorelasi dengan otentisitas (Mohamed \& Sarwar, 2022; Shaaban et al., 2026). Pemodelan gaya profetik dengan \textit{AraBERT} memperkuat identifikasi ini melalui sidik stilistik profetik sebagai penakar deviasi (Shaaban et al., 2026; Haque et al., 2020).

Pendekatan rekomendasi bukanlah parameter tunggal, melainkan sintesis kumulatif dan kontekstual antara \textit{sanad}, \textit{matan}, dan latar historis (Muhamed Ali et al., 2015; Wasman et al., 2023). Rantai yang lemah hanya menetapkan status \textit{dhaif}; vonis fabrikasi yang meyakinkan memerlukan konvergensi bukti berupa transmiter fraudulen, kontradiksi tekstual, absennya dukungan historis-sumber, transmisi yang mustahil, dan pengakuan pelaku (Wasman et al., 2023; Calgan, 2024). Kritik kontekstual menekankan bahwa \textit{sabab wurud} turut perlu diuji, sebab telaah \textit{sanad}-\textit{matan} saja tidak selalu menyingkap fabrikasi yang dipoles rapi (Wasman et al., 2023; Noor, 2023). Perbandingan klasifier \textit{deep learning}, kompresi, dan \textit{machine learning} tradisional memperlihatkan keunggulan kombinasi fitur struktural-semantik dibanding pendekatan tunggal --- selaras dengan kaidah kumulatif klasik (Tarmom et al., 2022; Haque et al., 2020). Tinjauan komputasional-NLP mengonfirmasi bahwa prasarana semantik modern meneguhkan nalar klasik: autentikasi menuntut multi-metode bagi setiap laporan (Azmi et al., 2019; Kamran et al., 2023). Penguasaan ciri \textit{sanad}-\textit{matan} mengharuskan institusionalisasi respons, sebab identifikasi personal tidak mencukupi untuk menghadapi industri yang sistematis (Muhamed Ali et al., 2015; Noor, 2023). Identifikasi otomatis melalui analisis sentimen dan \textit{machine learning} hanya signifikan bila dilatih dengan taksonomi klasik --- pembaruan algoritmik tetap bertumpu pada khazanah kriteria ulama (Haque et al., 2020; Tarmom et al., 2022). Kebangkitan tradisi Mamluk melalui \textit{Muqaddimah} karya Ibn al-Salah membuktikan bahwa setiap krisis direspons melalui kodifikasi metodologis yang segar, sebuah pola repetitif dari era klasik hingga digital (Noor, 2023; Su, 2021). Telaah Ibn Kasir terhadap periwayatan kanonik memperlihatkan kultur kritik yang tanpa area imun: keterbukaan pada pengujian ulang itulah yang menjadi proteksinya (Calgan, 2024; Wasman et al., 2023). Berangkat dari kebutuhan institusionalisasi tersebut, uraian berpindah menuju arsitektur penyelamatan: kerangka metodologis bagi pemurnian korpus secara permanen (Muhamed Ali et al., 2015; Noor, 2023).

\subsubsection*{5. Upaya Penyelamatan Hadis}

Upaya pertama ialah \textit{takhrij}: penelusuran asal genetik sekaligus verifikasi terpadu antara \textit{sanad} dan \textit{matan} (Muhamed Ali et al., 2015; Wasman et al., 2023). Seluruh jalur suatu \textit{matan} ditelusuri oleh peneliti, varian antarjalur dikomparasikan, kredibilitas tiap transmiter dinilai, dan koherensinya terhadap Qur'an, hadis sahih, serta fakta historis diuji sebelum vonis dijatuhkan (Wasman et al., 2023; Su, 2021). Tahapan \textit{takhrij} kini dibakukan pula sebagai fitur komputasional --- asesmen linguistik berbasis \textit{machine learning} dan prediksi sentimen melipatgandakan kecepatan penelusuran lintas-korpus tanpa meninggalkan nalar klasik (Mohamed \& Sarwar, 2022; Haque et al., 2020). Ontologi \textit{SemanticHadith} mengotomatisasi proses ini: korpus direpresentasikan sebagai jaringan semantik yang tertelusur, terbandingkan, dan teraudit secara sistematis (Kamran et al., 2023; Azmi et al., 2019). Dengan demikian, \textit{takhrij} berfungsi sebagai prosedur operasional yang menerjemahkan kriteria teoretis tentang ciri \textit{maudu'} menjadi putusan praktis bagi setiap riwayat yang bersirkulasi (Muhamed Ali et al., 2015; Wasman et al., 2023). Upaya kedua ialah \textit{al-jarh wa al-ta'dil}: penaksiran \textit{'adalah} dan \textit{dabt} setiap mata rantai sebagai benteng epistemologis (Muhamed Ali et al., 2015; Su, 2021). Disiplin ini semata-mata mengesahkan transmisi dari transmiter \textit{thiqah} sebagai hujah, sembari memfiltrasi sosok \textit{majhul}, pelupa berat, dan pendusta sebelum mencemari korpus normatif (Muhamed Ali et al., 2015; Noor, 2023). Telaah perintis 'Ali b. al-Madini meletakkan dasar bagi \textit{'ilal} dan \textit{rijal} yang kemudian diadopsi oleh Ibn al-Salah dan era Mamluk sebagai kurikulum baku (Su, 2021; Noor, 2023). Analisis jaringan transmiter \textit{Sahih Bukhari} memvalidasinya secara kuantitatif: transmiter poros yang tertaut padat dinilai lebih andal dibanding simpul pinggiran yang menyimpang (Alam \& Schneider, 2020; Kamran et al., 2023). Melalui \textit{jarh-ta'dil}, komunitas merancang pertahanan bertingkat yang terekam dalam literatur tersendiri (Muhamed Ali et al., 2015; Wasman et al., 2023).

Upaya ketiga ialah kodifikasi \textit{maudu'at}: penghimpunan laporan-laporan palsu menjadi arsip anti-disinformasi (Noor, 2023; Calgan, 2024). Kontribusi paling tegas disumbangkan oleh Ibn al-Jauzi melalui \textit{al-Maudhu'at}, yang menghimpun sekitar 1.600 laporan palsu --- sebelumnya terserak di berbagai literatur \textit{rijal} --- menjadi kritik tematik yang sistematis (Noor, 2023; Wasman et al., 2023). Tradisi ini dilanjutkan oleh al-Suyuti (\textit{al-La'ali al-Masnu'ah}), Ibn 'Iraq (\textit{Tanzih al-Syari'ah}), Syaukani (\textit{al-Fawa'id al-Majmu'ah}), dan pada masa modern oleh al-Albani (\textit{Silsilah al-Ahadits al-Dha'ifah wa al-Maudhu'ah}) yang menghidupkan kembali sensibilitas verifikasi revivalis (Noor, 2023; Muhamed Ali et al., 2015). Kritik Ibn Kasir yang mengoreksi sebagian materi \textit{Sahihayn} membuktikan bahwa pembakuan kritis mencakup kanon tanpa kecuali (Calgan, 2024; Wasman et al., 2023). Korpus \textit{maudu'at} klasik kini memiliki dwifungsi sebagai bahan latih model komputasional --- khazanah klasik menopang pembaruan digital (Tarmom et al., 2022; Shaaban et al., 2026).

Sebagai rangkuman Bab II: fabrikasi bukanlah deviasi pinggiran, melainkan krisis teologis kolosal yang mendesak peradaban Islam menyusun pertahanan epistemologis paling mutakhir --- terbentang dari \textit{tathabbut} para sahabat hingga graf pengetahuan digital masa kini (Muhamed Ali et al., 2015; Kamran et al., 2023). Survei komputasional, asesmen linguistik otomatis, dan pemodelan \textit{AraBERT} menunjukkan bahwa pembaruan algoritmik berperan sebagai penerus etos \textit{takhrij} dan \textit{jarh-ta'dil}, bukan sebagai substitusi bagi ulama (Azmi et al., 2019; Mohamed \& Sarwar, 2022; Shaaban et al., 2026). Akan tetapi, digitalisasi turut mengakselerasi peredaran kutipan yang belum terverifikasi, sehingga dependensi terhadap penilaian tradisional \textit{isnad}-\textit{matan} justru semakin menguat (Azmi et al., 2019; Wasman et al., 2023). Setelah periodisasi, penghimpunan, pemalsuan, dan penyelamatan dipetakan, uraian ini siap disintesiskan dalam Bab III: kesimpulan integratif dan agenda verifikasi digital yang berlandaskan pada kritik klasik (Noor, 2023; Kamran et al., 2023).

